%%%%%%%%%%%%%%%%%%%%%%%%%%%%%%%%%%%%%%%%%%%%%%%%%%%%%%%%%%%%%%%%%%%%%%%%%%%%%%%%
%%  EINLAGE 1: ARBEITSPLAN          Zwischenabgabe: 28.10.2026 (obligatorisch)
%%  Struktur: Beschreibung (1/4–1/2 S.) + Reflexion (1–2 S.)
%%  Die Einlage selbst (der Plan) steht im Anhang: 1_arbeitsplan_anhang.tex
%%%%%%%%%%%%%%%%%%%%%%%%%%%%%%%%%%%%%%%%%%%%%%%%%%%%%%%%%%%%%%%%%%%%%%%%%%%%%%%%
\section{Einlage: Arbeitsplan}
\label{sec:arbeitsplan}

\begin{hinweise}
Diese Einlage ist Ihr Arbeitsplan. Der Arbeitsplan ist die \emph{einzige} Einlage, die mit der Partnerin/dem
Partner identisch sein darf (\enquote{Ihr gemeinsamer Arbeitsplan}). Hinweise zur Erstellung (Marti): nicht nur
eine grobe Liste, sondern feine Teilaufgaben mit zeitlichem Ablauf, sichtbar machen, was parallel läuft;
Vorbereitung (\enquote{Einkaufen}) und Nachbereitung (\enquote{Abwasch}) nicht vergessen -- z.\,B. als
Excel-Tabelle oder in ProjectLibre, oder mit dem Gantt-Gerüst im Anhang.
\end{hinweise}

\subsection{Beschreibung}
\begin{hinweise}
Die Beschreibung gibt alle notwendigen Informationen an, die die Einlage beschreiben, also zum Beispiel den
Fundort, die Literaturstelle und eine kurze Inhaltsübersicht, aber keine Nacherzählung ($1/4$ Seite, max.
$1/2$ Seite!). Leitfaden: Um welche Art von Einlage handelt es sich? Woher stammt sie? Welches sind die
wesentlichen Aussagen?
\end{hinweise}
\begin{umfang}{Einlage 1 (Arbeitsplan): Beschreibung}{0.25}{0.5}
\platzhalter{Art der Einlage (z.\,B. Gantt-Diagramm, erstellt mit \dots), wann/wie gemeinsam mit \Partner{}
erstellt, Zeitraum, Hauptphasen, Meilensteine; Verweis auf Anhang~\ref{anh:arbeitsplan}}
\end{umfang}

\subsection{Reflexion}
\begin{hinweise}
In der Reflexion sollen Sie sich Gedanken machen über die Wirkung dieser Einlage und der Arbeit mit ihr auf
Ihren Lernfortschritt.
\begin{itemize}\itemsep0pt
\item In welchem Zusammenhang steht der Arbeitsplan zu Ihrem Projektthema?
\item In welchem Zusammenhang steht der Arbeitsplan zu Ihrem Schwerpunktthema?
\item Was haben Sie bei der Erstellung des Arbeitsplanes gelernt?
\item Wie hat Ihnen der Arbeitsplan beim Verständnis Ihres Projektthemas und Ihres Schwerpunktthemas geholfen?
\item Wo gab es Schwierigkeiten bei der Beschäftigung mit der Arbeitsplan?
\item Was sagen Ihre Kommilitoninnen oder Kommilitonen dazu, als Kommentar oder wenn Sie mit Ihnen über den Arbeitsplan diskutieren?
\item Enthält Ihr Arbeitsplan logische oder physikalische Fehler?
\item Gab es bei der Beschäftigung mit dem Arbeitsplan eine überraschende neue Erkenntnis?
\end{itemize}
\textbf{Die Reflexion für eine Einlage sollte mindestens eine Seite und maximal 2 Seiten lang sein. Die
Längenangaben sind ohne Titel, Bilder, Tabellen und Tabellen- oder Bildlegenden zu verstehen.}
Peer-Feedback zitieren (siehe \texttt{literatur.bib}, Vorlage \texttt{VORLAGE-Gespraech}).
\end{hinweise}
\begin{umfang}{Einlage 1 (Arbeitsplan): Reflexion}{1}{2}
\platzhalter{Reflexion in Fließtext -- Leitfragen erkennbar beantworten}
\end{umfang}

%% Optional (vom Leitfaden ausdrücklich erwünscht): eigene Bewertungskriterien
%\subsection{Eigene Bewertungskriterien}
%Für diese Einlage schlage ich zusätzlich folgende Kriterien vor: ...
